\section{固定制度与稳态对象下的市场出清条件}





















































































































































































































































































































本阶段在给定 candidate incumbent measures 与 entrant masses 的条件下，推导 benchmark 的市场出清模块。目标是把合格受托生产市场与劳动力市场出清条件写成可验证、可求解的形式。这里不引入任何新机制，只把使均衡方程良定义并可求解的正则条件显式写出。这里使用的 incumbent measures 只是价格出清模块的给定输入；Phase 7 才会进一步加入 stationary consistency。

\subsection{Benchmark Fidelity Check}

\paragraph{Step 1：本阶段使用的 benchmark 对象。}
本阶段的 canonical 对象是：
\begin{itemize}
    \item 来自 type-$B$ 静态模块的 $x_B^*(z;p_m,M)$ 与 $\lambda_B(z,p_m,M)$；
    \item 来自 type-$C$ 静态模块的 $m^*(z;p_m,w)$；
    \item 本阶段视为给定的 candidate incumbent measures $\mu_A,\mu_B,\mu_C$，以及 entrant masses $n_E^A,n_E^B,n_E^C$；
    \item 分类型劳动需求对象 $\ell_A(z;w)$、$\ell_B(z;p_m,M)$、$\ell_C(z;p_m,w)$，以及进入者启动劳动系数 $\ell_E^j$；
    \item 外生劳动供给 $\bar L$。
\end{itemize}

\paragraph{Step 2：benchmark 方程。}
合格受托生产市场的需求与供给分别定义为
\begin{align}
D_m(p_m,M)
&=\int_Z x_B^*(z;p_m,M)\,\lambda_B(z,p_m,M)\,\mu_B(dz),
\label{eq:p6zh_Dm}\\
S_m(p_m,w)
&=\int_Z m^*(z;p_m,w)\,\mu_C(dz),
\label{eq:p6zh_Sm}
\end{align}
市场出清条件是
\begin{equation}
D_m(p_m,M)=S_m(p_m,w).
\label{eq:p6zh_pm_clear}
\end{equation}

主文后来把 labor block 压缩写成 equilibrium object $L_D(M)$。为了避免静默扩张 benchmark 的 argument list，本阶段引入如下辅助的 fixed-price labor-demand map：
\begin{equation}
\begin{aligned}
\mathfrak L_D(w,p_m,M)={}&\int_Z \ell_A(z;w)\,\mu_A(dz)
+\int_Z \ell_B(z;p_m,M)\,\mu_B(dz)\\
&+\int_Z \ell_C(z;p_m,w)\,\mu_C(dz)
+\sum_{j\in\{A,B,C\}}\ell_E^j n_E^j(M),
\end{aligned}
\label{eq:p6zh_LD_aux}
\end{equation}
劳动市场出清条件是
\begin{equation}
\mathfrak L_D(w,p_m,M)=\bar L.
\label{eq:p6zh_labor_clear}
\end{equation}

\paragraph{Step 3：benchmark-fidelity 边界。}
contract-manufacturing block 可以直接由上面的 benchmark objects 写出。labor block 也可以忠实写出，但固定价格记号 $\mathfrak L_D(w,p_m,M)$ 只是为了避免把 benchmark object $L_D(M)$ 静默改写成扩展 argument list 的辅助表示。除此之外，当前 benchmark sources 没有给出足够细的 primitive labor-demand 推导，使我们能够从更低层对象自动推出 aggregate labor-demand continuity，所以后文的 labor 正则条件会明确写成补充闭合条件。

\paragraph{Step 4：边界纪律。}
本阶段严格保持 benchmark 边界：不加入政府最优化、不做 DSGE 化、不引入并购/议价层，也不扩展核心状态维度。

\subsection{Part 1：合格受托生产市场}

\paragraph{需求与供给良定义所需条件。}
固定制度 $M$ 与工资 $w$。
\begin{assumption}[MC1：可积性]
对每个可行的 $p_m$，映射 $z\mapsto x_B^*(z;p_m,M)\lambda_B(z,p_m,M)$ 与 $z\mapsto m^*(z;p_m,w)$ 在 Borel 意义下可测，并且分别对 $\mu_B,\mu_C$ 可积。
\end{assumption}

在写 aggregate continuity assumption 之前，先把前序阶段已经给出的内容与本阶段新增的条件分开。Phase 2 已经推导出 $x_B^*(z;p_m,M)$ 与 $m^*(z;p_m,w)$ 在紧价格域上的连续性。因此，只要 $\lambda_B(z,p_m,M)$ 关于 $p_m$ 连续且有界，并且候选测度 $\mu_B,\mu_C$ 有限，就可以通过 dominated-convergence argument 推出下面 aggregate demand/supply maps 的连续性。MC2 在本阶段把这一 aggregate regularity 明确记成条件，因为市场出清定理真正使用的是 aggregate objects，而不是逐点 policy functions。

\begin{assumption}[MC2：关于 $p_m$ 的连续性]
$D_m(p_m,M)$ 与 $S_m(p_m,w)$ 在紧价格区间 $\mathcal P_m=[\underline p_m,\overline p_m]\subset\mathbb R_+$ 上关于 $p_m$ 连续。
\end{assumption}

\begin{assumption}[MC3：单调性]
在 $\mathcal P_m$ 上，$D_m(\cdot,M)$ 弱递减，$S_m(\cdot,w)$ 弱递增。
\end{assumption}

\begin{assumption}[MC4：端点跨越条件]
在区间端点满足
\[
D_m(\underline p_m,M)-S_m(\underline p_m,w)\ge 0,
\qquad
D_m(\overline p_m,M)-S_m(\overline p_m,w)\le 0.
\]
\end{assumption}

\begin{proposition}[合格受托生产市场出清价格的存在性]
在 MC1--MC4 下，至少存在一个 $p_m^*(w,M)\in\mathcal P_m$ 使得
\[
D_m(p_m^*,M)=S_m(p_m^*,w).
\]
若进一步假设 $D_m(\cdot,M)$ 严格递减且 $S_m(\cdot,w)$ 严格递增，则该出清价格唯一。
\end{proposition}

\begin{proof}
定义超额需求
\[
\Delta_m(p_m;w,M):=D_m(p_m,M)-S_m(p_m,w).
\]
由 MC2，$\Delta_m$ 在紧区间 $\mathcal P_m$ 上连续。由 MC4，
\[
\Delta_m(\underline p_m;w,M)\ge 0,
\qquad
\Delta_m(\overline p_m;w,M)\le 0.
\]
据介值定理，至少存在一个 $p_m^*$ 使 $\Delta_m(p_m^*;w,M)=0$。

若再有严格单调性，设 $p_1<p_2$ 都是出清解，则应有
\[
\Delta_m(p_1;w,M)>\Delta_m(p_2;w,M),
\]
与两者都等于 0 矛盾，故解唯一。
\end{proof}

\subsection{Part 2：劳动力市场}

\paragraph{劳动需求闭合所需条件。}
固定制度 $M$ 与候选价格 $p_m$。
与 contract-manufacturing block 不同，当前 benchmark sources 并没有给出足够细的 primitive labor-demand 推导，使我们可以从更低层对象自动推出 aggregate labor-demand continuity。因此，下面的条件最好理解为施加在辅助 fixed-price labor-demand map 上的补充闭合条件。
\begin{assumption}[LC1：可积性与连续性]
\eqref{eq:p6zh_LD_aux} 定义的 $\mathfrak L_D(w,p_m,M)$ 在紧工资区间 $\mathcal W=[\underline w,\overline w]\subset\mathbb R_+$ 上关于 $w$ 有限且连续。
\end{assumption}

\begin{assumption}[LC2：关于工资的单调性]
$\mathfrak L_D(\cdot,p_m,M)$ 在 $\mathcal W$ 上弱递减。
\end{assumption}

\begin{assumption}[LC3：端点跨越条件]
\[
\mathfrak L_D(\underline w,p_m,M)\ge \bar L,
\qquad
\mathfrak L_D(\overline w,p_m,M)\le \bar L.
\]
\end{assumption}

\begin{proposition}[劳动力市场出清工资的存在性]
在 LC1--LC3 下，对每个固定的 $(p_m,M)$，至少存在一个 $w^*(p_m,M)\in\mathcal W$ 使
\[
\mathfrak L_D(w^*,p_m,M)=\bar L.
\]
若 $\mathfrak L_D(\cdot,p_m,M)$ 严格递减，则该出清工资唯一。
\end{proposition}

\begin{proof}
定义超额劳动需求
\[
\Delta_L(w;p_m,M):=\mathfrak L_D(w,p_m,M)-\bar L.
\]
由 LC1，$\Delta_L$ 在 $\mathcal W$ 上连续。由 LC3，
\[
\Delta_L(\underline w;p_m,M)\ge 0,
\qquad
\Delta_L(\overline w;p_m,M)\le 0.
\]
据介值定理，至少存在一个根 $w^*$。

若 $\mathfrak L_D(\cdot,p_m,M)$ 严格递减，则 $\Delta_L(\cdot;p_m,M)$ 也严格递减，只能穿过 0 一次，所以解唯一。
\end{proof}

\subsection{Part 3：联合价格出清映射（接口声明）}

两个市场通过 $(p_m,w)$ 联立耦合。本阶段已建立“给定另一个价格时”的一维存在性结果。下一阶段可以把
\[
\begin{cases}
D_m(p_m,M)-S_m(p_m,w)=0,\\
\mathfrak L_D(w,p_m,M)-\bar L=0,
\end{cases}
\]
作为定义在 $\mathcal P_m\times\mathcal W$ 上的联合不动点系统处理。

\paragraph{为什么这已满足 Phase 6 目标。}
Phase 6 不需要重解 Bellman，也不需要重做 entry 推导；它要完成的是把 market-clearing closure 写成规范且可验证的接口，供后续稳态分布与完整均衡存在性阶段调用。

\subsection{What did we just do, and why does it matter?}

本阶段把连接企业动态问题与后续均衡闭合的两个价格出清条件单独拆开处理。整个推导过程中，incumbent measures 被当作候选输入，而不是已经先于 Phase 6 推导完成的稳态对象；Phase 7 才负责加入 stationary consistency。对 contract-manufacturing 市场而言，核心是一维方程
\[
D_m(p_m,M)=S_m(p_m,w),
\]
它要求 benchmark 下对合格受托生产服务的需求，必须等于在候选价格下由 type-$C$ incumbents 提供的供给。对 labor market 而言，核心是一维方程
\[
\mathfrak L_D(w,p_m,M)=\bar L,
\]
它是 benchmark labor-clearing object $L_D(M)=\bar L$ 的 fixed-price auxiliary representation。把每个市场出清条件分别写成一维问题并单独证明存在性，有两个好处：第一，连续性把问题变成良定义的 root-finding problem；第二，单调性与端点跨越条件把介值定理真正能够使用的地方逐项写清楚。

这一步之所以重要，是因为后面的均衡存在性论证不能只说“价格会存在”。它必须明确：到底是哪两个标量方程在约束价格，这些方程依赖哪些候选对象，以及这些条件如何被放进后续的有限维 outer map。Phase 6 的作用正是在这里。它不重写企业优化问题，不额外添加状态或控制，也不提前宣称 full equilibrium 已经存在；它只是把 benchmark 的 market-clearing block 写成一个可以被下一阶段直接调用的、定理级别可核对的接口。

\subsection*{End-of-Section Recap}

\begin{itemize}
    \item What has been established mathematically: 合格受托生产市场与劳动力市场的出清方程都是良定义的，并且在显式的连续性、单调性与端点跨越条件下分别有存在性结论。
    \item What assumptions were used: 本阶段视为给定的 candidate incumbent measures 与 entrant masses，以及 contract-manufacturing block 的 MC1--MC4 和辅助 fixed-price labor block 的 LC1--LC3。
    \item What remains to be derived next: 下一阶段必须对这些 candidate incumbent measures 加入 stationary consistency，并把所得稳态对象与市场出清条件合并起来。
\end{itemize}
